% Generated by scripts/assemble.py; corrections belong in source/corrections.json.
\examhead{68}{太原理工大学 2026 年数学分析试卷}{68}
% p069-b07; source PDF page 69
\par\Needspace{6\baselineskip}
\blockanchor{p069-b07}
\question{1}求极限\par
\nopagebreak[4]
% p069-b08; source PDF page 69
\blockanchor{p069-b08}
\displayformula{\lim _ {n \rightarrow \infty} \left(\frac {\sin \frac {\pi}{n}}{n + 2026} + \frac {\sin \frac {2 \pi}{n}}{n + \frac {2026}{2}} + \dots + \frac {\sin \frac {n \pi}{n}}{n + \frac {2026}{n}}\right)}
% p069-b09; source PDF page 69
\blockanchor{p069-b09}
\question{2}求极限 $\lim_{x\to 0}\frac{\tan(\tan x) - \sin(\sin x)}{x^2\arctan 2x}.$\par
% p069-b10; source PDF page 69
\blockanchor{p069-b10}
\question{3}计算不定积分 $\int \frac{\cos^2 x}{(a^2 \sin^2 x + b^2 \cos^2 x)^2} \, \mathrm{d}x$.\par
% p069-b11; source PDF page 69
\blockanchor{p069-b11}
\question{4}计算曲线  $\left(\frac{x^{2}}{a^{2}}+\frac{y^{2}}{b^{2}}\right)^{2}=x^{2}+y^{2}(a,b>0)$  围成的平面图形的面积.\par
% p069-b12; source PDF page 69
\par\Needspace{6\baselineskip}
\blockanchor{p069-b12}
\question{5}计算曲线积分\par
\nopagebreak[4]
% p069-b13; source PDF page 69
\blockanchor{p069-b13}
\displayformula{\oint_ {L} (z - y) \mathrm{d} x + (x - z) \mathrm{d} y + (x - y) \mathrm{d} z.}
% p069-b14; source PDF page 69
\blockanchor{p069-b14}
其中 L 是  $x^{2} + y^{2} = 4$  与  $x - y + z = 1$  的交线，从 z 轴正向往 z 轴负向看是顺时针方向.\par
% p069-b15; source PDF page 69
\par\Needspace{6\baselineskip}
\blockanchor{p069-b15}
\question{6}计算第二型曲面积分\par
\nopagebreak[4]
% p069-b16; source PDF page 69
\blockanchor{p069-b16}
\displayformula{\iint_ {S} 2 (1 - x ^ {2}) \mathrm{d} y \mathrm{d} z + 8 x y \mathrm{d} z \mathrm{d} x - 4 x z \mathrm{d} x \mathrm{d} y.}
% p069-b17; source PDF page 69
\blockanchor{p069-b17}
其中 $S$ 是 $x = 5^{y}$$(0\leq y\leq 1)$ 绕 $x$ 轴形成的曲面的外侧.\par
% p069-b18; source PDF page 69
\blockanchor{p069-b18}
\question{7}计算 $\int_0^{+\infty} \frac{\arctan x}{x(1 + x^2)} \, \mathrm{d}x$.\par
% p069-b19; source PDF page 69
\blockanchor{p069-b19}
\question{8}判断反常积分 $\int_0^{+\infty} \frac{x \ln x}{x^p + 1} (p > 0)$ 的敛散性.\par
% p069-b20; source PDF page 69
\blockanchor{p069-b20}
\question{9}设 $f(x)$ 是 $2\pi$ 为周期的函数, $f(x) = x(\pi - x)$, $x \in [0, \pi]$.\par
% p069-b21; source PDF page 69
\blockanchor{p069-b21}
(1) 求 $f(x)$ 的余弦级数和正弦级数展开式.\par
% p069-b22; source PDF page 69
\blockanchor{p069-b22}
(2) 求 $\sum_{n=1}^{\infty} \frac{1}{n^2}$ 和 $\sum_{n=1}^{\infty} \frac{(-1)^{n-1}}{(2n-1)^3}$.\par
% p069-b23; source PDF page 69
\par\Needspace{6\baselineskip}
\blockanchor{p069-b23}
\question{10}设函数 $f(x)$ 在 $[-1, 1]$ 上连续，令 $\varphi_n(x) = \begin{cases} (1 - x)^n, & 0 \leq x \leq 1; \\ (1 + x)^n, & -1 \leq x \leq 0. \end{cases}$ 证明：\par
\nopagebreak[4]
% p069-b24; source PDF page 69
\blockanchor{p069-b24}
\displayformula{\lim _ {n \rightarrow \infty} \frac {n + 1}{2} \int_ {- 1} ^ {1} f (x) \varphi_ {n} (x) \mathrm{d} x = f (0).}
% p069-b25; source PDF page 69
\blockanchor{p069-b25}
\question{11}数列 $\{a_{n}\}$ 满足 $a_0 = 1, a_1 = 1, a_n = a_{n-1} + 2a_{n-2} (n \geq 2)$.\par
% p069-b26; source PDF page 69
\blockanchor{p069-b26}
(1) 证明: $\left(\frac{5}{3}\right)^{n-1} \leq a_n \leq 3^{n-1}$.\par
% p070-b00; source PDF page 70
\blockanchor{p070-b00}
(2) 证明: $\sum_{n=1}^{\infty} \frac{1}{\ln a_n}$ 发散, $\sum_{n=1}^{\infty} \frac{1}{a_n}$ 收敛.\par
% p070-b01; source PDF page 70
\blockanchor{p070-b01}
(3) 若 $\lim_{n\to \infty}\frac{a_n}{a_{n + 1}}$ 存在, 求幂级数 $\sum_{n = 1}^{\infty}a_nx^n$ 的收敛半径以及收敛区间内的和函数.\par
% p070-b02; source PDF page 70
\blockanchor{p070-b02}
\question{12}解答如下问题:\par
% p070-b03; source PDF page 70
\blockanchor{p070-b03}
(1) 设 $f(x), g(x)$ 都是凸函数, $f(x)g(x)$ 和 $f(g(x))$ 是否是凸函数? 若不是, 给出反例. 对于 $f(g(x))$, 如何加一个充分条件, 使之成为凸函数? 给出证明.\par
% p070-b04; source PDF page 70
\blockanchor{p070-b04}
(2) 证明: 实数域上的有界凸函数必是常值函数.\par
